\section{Source-only reduction of distributed densities}
\label{sec:peps_source_only_reduction}

We prove the source-only reduction from
\cite[proof of Theorem~5.2]{OpenAI2026PolynomialPEPS}.
% Source: 04-compression.tex, lines 199--229, at
% adc7f1241b42e322a6451854ab7e4b4c146bf78a.
All registers are retained
until the final partial trace. The Hilbert spaces have arbitrary finite
dimensions.


\begin{definition}[Private product inputs and exact local operations]
    \label{def:peps_private_product_input}
    \lean{TNLean.PEPS.PairEffect.ProductInput.mem_owners_wholePartyLayout}
    \lean{TNLean.PEPS.PairEffect.ProductInput.nodup_owners_wholePartyLayout}
    \lean{TNLean.PEPS.PairEffect.ProductInput.ofAllParties}
    \lean{TNLean.PEPS.PairEffect.ProductInput.owners_wholePartyLayout}
    \lean{TNLean.PEPS.PairEffect.ProductInput.wholePartyLayout}
    \lean{TNLean.PEPS.PairEffect.ProductInput}
    \lean{TNLean.PEPS.PairEffect.ProductInput.eval_prepare}
    \lean{TNLean.PEPS.PairEffect.ProductInput.eval_prepare_smul}
    \lean{TNLean.PEPS.PairEffect.ProductInput.isAllowed_prepare}
    \lean{TNLean.PEPS.PairEffect.ProductInput.mapOwner}
    \lean{TNLean.PEPS.PairEffect.ProductInput.mapOwnerIso_vector}
    \lean{TNLean.PEPS.PairEffect.ProductInput.norm_vector}
    \lean{TNLean.PEPS.PairEffect.ProductInput.ofParties}
    \lean{TNLean.PEPS.PairEffect.ProductInput.partitionIso_symm_mapOwner_vector}
    \lean{TNLean.PEPS.PairEffect.ProductInput.partitionIso_symm_vector}
    \lean{TNLean.PEPS.PairEffect.ProductInput.partitionIso_vector}
    \lean{TNLean.PEPS.PairEffect.ProductInput.prepare}
    \lean{TNLean.PEPS.PairEffect.ProductInput.prepareIsometry}
    \lean{TNLean.PEPS.PairEffect.ProductInput.restrict}
    \lean{TNLean.PEPS.PairEffect.ProductInput.sources_prepare}
    \lean{TNLean.PEPS.PairEffect.ProductInput.vector}
    \leanok
    \uses{def:peps_party_layout,def:peps_party_monomial}
    For an ordered register list with spaces $H_1,\ldots,H_n$, choose
    unit vectors $x_j\in H_j$ and put $x=x_1\otimes\cdots\otimes x_n$.
    The empty tensor is the scalar $1$. Registers may instead be grouped
    into one entire private memory for each party, so no factorization
    inside a party's memory is required. Local preparation is the
    isometry $z\mapsto zx$ from $\C$ to the input memory.

    Restricting to a set of owners selects the corresponding factors.
    The canonical owner partition sends $x$ to $x_A\otimes x_{A^c}$;
    each factor has norm one. Changing owner names leaves these vectors
    unchanged under the canonical memory identification.
\end{definition}

\begin{definition}[Prepared sources and retained spectators]
    \label{def:peps_source_replacement_operations}
    \lean{TNLean.PEPS.PairEffect.EffectCircuit}
    \lean{TNLean.PEPS.PairEffect.EffectCircuit.IsAllowed}
    \lean{TNLean.PEPS.PairEffect.EffectCircuit.eval}
    \lean{TNLean.PEPS.PairEffect.EffectCircuit.norm_eval_le_one}
    \lean{TNLean.PEPS.PairEffect.PartyChain.isPreparation_prepStack}
    \lean{TNLean.PEPS.PairEffect.SourceCircuit.castLayouts}
    \lean{TNLean.PEPS.PairEffect.SourceCircuit.eval_castLayouts}
    \lean{TNLean.PEPS.PairEffect.SourceCircuit.eval_exchangeBlocks}
    \lean{TNLean.PEPS.PairEffect.SourceCircuit.eval_frameList_tmul}
    \lean{TNLean.PEPS.PairEffect.SourceCircuit.eval_ofLocalWord}
    \lean{TNLean.PEPS.PairEffect.SourceCircuit.exchangeBlocks}
    \lean{TNLean.PEPS.PairEffect.SourceCircuit.frameList}
    \lean{TNLean.PEPS.PairEffect.SourceCircuit.isAllowed_castLayouts}
    \lean{TNLean.PEPS.PairEffect.SourceCircuit.isAllowed_frameList}
    \lean{TNLean.PEPS.PairEffect.SourceCircuit.isAllowed_ofLocalWord}
    \lean{TNLean.PEPS.PairEffect.SourceCircuit.ofLocalWord}
    \lean{TNLean.PEPS.PairEffect.SourceInventory.norm_vector}
    \lean{TNLean.PEPS.PairEffect.Word.IsPreparation}
    \lean{TNLean.PEPS.PairEffect.Word.eval_isPreparation_cast}
    \lean{TNLean.PEPS.PairEffect.Word.eval_isPreparation_heq}
    \lean{TNLean.PEPS.PairEffect.Word.eval_localPrepare}
    \lean{TNLean.PEPS.PairEffect.Word.isAllowed_localPrepare}
    \lean{TNLean.PEPS.PairEffect.Word.localPrepare}
    \lean{TNLean.PEPS.PairEffect.Word.output_of_isPreparation}
    \lean{TNLean.PEPS.PairEffect.exists_prepared_effectReplacement}
    \lean{TNLean.PEPS.PairEffect.isPreparation_prepGate}
    \lean{TNLean.PEPS.PairEffect.isPreparation_prepWord}
    \leanok
    \uses{def:peps_private_product_input,def:peps_party_monomial}
    A source-only operation is composed from local contractions,
    register permutations, and normalized pair-source preparations.
    A local vector $x$ is prepared by $y\mapsto x\otimes y$.
    Preparing an ordered inventory of pair vectors gives the tensor
    product of those vectors followed by the unchanged input.
    Its norm is one when every pair vector has norm one.

    For an operation $T$ and an additional register memory $R$, retaining
    $R$ means applying $I_R\otimes T$. Exchanging two register blocks
    sends $r\otimes s\otimes y$ to $s\otimes r\otimes y$.
    Equal register lists are identified by the canonical memory isometry.
    These operations preserve the indicated vectors and introduce no
    additional nonprivate gate. An allowed original monomial expansion
    can consequently be replaced by its actual prepared-source expansion.
\end{definition}

\begin{definition}[Original distributed circuits]
    \label{def:peps_original_circuit}
    \lean{TNLean.PEPS.PairEffect.EffectCircuit.gateLocations}
    \lean{TNLean.PEPS.PairEffect.EffectCircuit.participants}
    \lean{TNLean.PEPS.PairEffect.OriginalCircuit.isAllowed_produce}
    \lean{TNLean.PEPS.PairEffect.OriginalCircuit.isExpansionBounded_produce_iff}
    \lean{TNLean.PEPS.PairEffect.OriginalCircuit.nonprivateLocations}
    \lean{TNLean.PEPS.PairEffect.OriginalCircuit.participants}
    \lean{TNLean.PEPS.PairEffect.OriginalCircuit.participants_produce}
    \lean{TNLean.PEPS.PairEffect.OriginalCircuit.produce}
    \lean{TNLean.PEPS.PairEffect.OriginalCircuit.produceLocationsEquiv}
    \lean{TNLean.PEPS.PairEffect.OriginalCircuit}
    \lean{TNLean.PEPS.PairEffect.OriginalCircuit.eval}
    \lean{TNLean.PEPS.PairEffect.OriginalCircuit.nonprivateCount}
    \lean{TNLean.PEPS.PairEffect.OriginalCircuit.IsExpansionBounded}
    \leanok
    \uses{def:peps_party_layout,def:peps_party_monomial}
    An original circuit is an ordered composition of contractions on
    owned registers. A private operation acts at one party. Every
    nonprivate operation has a specified participating set and an
    expansion $G_j=\sum_\xi c_{j,\xi}M_{j,\xi}$ into allowed monomials.
    Registers outside the participating set are spectators.
    Let $N$ denote the number of nonprivate occurrences. Fix $r$ and $S$
    so that each original monomial has at most $r$ pair effects and
    $\sum_\xi|c_{j,\xi}|\le S$ at every nonprivate occurrence.
\end{definition}

\begin{theorem}[Common auxiliary vector and operator error]
    \label{thm:peps_circuit_effect_error}
    \lean{TNLean.PEPS.PairEffect.EffectCircuit.referenceEval_apply}
    \lean{TNLean.PEPS.PairEffect.EffectCircuit.norm_auxiliaryVector}
    \lean{TNLean.PEPS.PairEffect.EffectCircuit.norm_replacement_sub_reference_le}
    \lean{TNLean.PEPS.PairEffect.EffectCircuit.isAllowed_replacement}
    \lean{TNLean.PEPS.PairEffect.OriginalCircuit.eval_produce}
    \lean{TNLean.PEPS.PairEffect.OriginalCircuit.expandedGateCount_produce}
    \leanok
    \uses{def:peps_original_circuit,cor:peps_pair_effects_rescaling}
    Let $T$ be the operator of an original circuit and let $\delta>0$.
    Choose $m=\lceil(rS/\delta)^2\rceil+1$, using the original expansions.
    Replacing each nonprivate operation by its rescaled source-only
    expansion gives a contraction $\widetilde T$ and a normalized
    auxiliary vector $\Gamma$ such that
    \begin{align}
        \|\widetilde T-|\Gamma\rangle\otimes T\|\le 2\delta N.
        \label{eq:peps_circuit_effect_error}
    \end{align}
    Each previously added auxiliary register is a spectator for all
    subsequent original operations.
\end{theorem}

\begin{proof}\leanok
    \lean{TNLean.PEPS.PairEffect.EffectCircuit.IsExpansionBounded}
    \lean{TNLean.PEPS.PairEffect.EffectCircuit.auxiliary}
    \lean{TNLean.PEPS.PairEffect.EffectCircuit.auxiliaryVector}
    \lean{TNLean.PEPS.PairEffect.EffectCircuit.expandedGateCount}
    \lean{TNLean.PEPS.PairEffect.EffectCircuit.replacement}
    \lean{TNLean.PEPS.PairEffect.SourceCircuit.eval_frameList}
    \lean{TNLean.PEPS.PairEffect.SourceCircuit.isAllowed_exchangeBlocks}
    \lean{TNLean.PEPS.PairEffect.appendIso_assoc_symm_heq}
    \lean{TNLean.PEPS.PairEffect.castOutputMap}
    \lean{TNLean.PEPS.PairEffect.castOutputMap_sub}
    \lean{TNLean.PEPS.PairEffect.exists_prepared_effectReplacement_native}
    \lean{TNLean.PEPS.PairEffect.exists_prepared_effectReplacement_uniform}
    \lean{TNLean.PEPS.PairEffect.gateMemoryMap}
    \lean{TNLean.PEPS.PairEffect.gateMemoryMap_sub}
    \lean{TNLean.PEPS.PairEffect.gate_auxiliary_layout}
    \lean{TNLean.PEPS.PairEffect.norm_castOutputMap}
    \lean{TNLean.PEPS.PairEffect.norm_gateMemoryMap_le}
    \lean{TNLean.PEPS.PairEffect.norm_retainMap_le}
    \lean{TNLean.PEPS.PairEffect.replacementCoefficients}
    \lean{TNLean.PEPS.PairEffect.retainMap}
    \lean{TNLean.PEPS.PairEffect.retainMap_sub}
    \lean{TNLean.PEPS.PairEffect.retainMap_tmul}
    \lean{TNLean.PEPS.PairEffect.EffectCircuit.norm_referenceEval_le_one}
    \lean{TNLean.PEPS.PairEffect.EffectCircuit.referenceEval}
    \lean{TNLean.PEPS.PairEffect.SourceCircuit.exchangeHead_tmul}
    \lean{TNLean.PEPS.PairEffect.SourceCircuit.norm_eval_castLayouts_sub}
    \lean{TNLean.PEPS.PairEffect.appendIso_singleton_symm_tmul}
    \lean{TNLean.PEPS.PairEffect.gateMemoryMap_auxiliary_identity}
    \lean{TNLean.PEPS.PairEffect.gateMemoryMap_auxiliary_tmul}
    \lean{TNLean.PEPS.PairEffect.gateMemoryMap_tmul}
    \lean{TNLean.PEPS.PairEffect.preparation_mapOwner_heq}
    \uses{cor:peps_pair_effects_rescaling}
    Induct on the composition of the original operations. The ideal
    replacement of each nonprivate gate is its original operator
    tensored with that gate's fixed normalized auxiliary vector.
    Their tensor product is $\Gamma$. Private contractions and
    permutations contribute no error. In a composition, insert the
    ideal intermediate operator and use contractivity to obtain
    $e(T_2T_1)\le e(T_1)+e(T_2)$. Each replaced nonprivate occurrence
    contributes at most $2\delta$, proving
    (\ref{eq:peps_circuit_effect_error}).
\end{proof}

\begin{theorem}[Preservation of participating parties]
    \label{thm:peps_replacement_participants}
    \lean{TNLean.PEPS.PairEffect.OriginalCircuit.replacementLocationsEquiv}
    \lean{TNLean.PEPS.PairEffect.OriginalCircuit.participants_replacementLocationsEquiv}
    \lean{TNLean.PEPS.PairEffect.OriginalCircuit.card_participants_ne_one}
    \leanok
    \uses{def:peps_original_circuit}
    The nonprivate occurrences of the original circuit correspond
    bijectively to the expanded gates of the replacement. Corresponding
    occurrences have exactly the same participating parties. Auxiliary
    preparations and register permutations introduce no additional
    expanded gate occurrences.
\end{theorem}

\begin{proof}\leanok
    \lean{TNLean.PEPS.PairEffect.EffectCircuit.exists_replacementLocationsEquiv}
    \lean{TNLean.PEPS.PairEffect.EffectCircuit.participants_replacementLocationsEquiv}
    \lean{TNLean.PEPS.PairEffect.EffectCircuit.replacementLocationsEquiv}
    \lean{TNLean.PEPS.PairEffect.SourceCircuit.castLocationsEquiv}
    \lean{TNLean.PEPS.PairEffect.SourceCircuit.exchangeBlocksLocationsEquivEmpty}
    \lean{TNLean.PEPS.PairEffect.SourceCircuit.false_of_gateLocation_ofLocalWord}
    \lean{TNLean.PEPS.PairEffect.SourceCircuit.frameListLocationsEquiv}
    \lean{TNLean.PEPS.PairEffect.SourceCircuit.ofLocalWordLocationsEquivEmpty}
    \lean{TNLean.PEPS.PairEffect.SourceCircuit.participants_castLocationsEquiv}
    \lean{TNLean.PEPS.PairEffect.SourceCircuit.participants_frameListLocationsEquiv}
    Induct on the circuit. An original nonprivate occurrence contributes
    one replacement gate with the original participating set. Composition
    joins the two occurrence sets, while spectator registers leave them
    unchanged. Private operations and permutations contribute empty sets.
\end{proof}

\begin{theorem}[Physical density error]
    \label{thm:peps_source_only_density_error}
    \lean{TNLean.PEPS.PairEffect.OriginalCircuit.rectangularTraceNorm_sourceOnly_density_sub_le}
    \lean{TNLean.PEPS.PairEffect.OriginalCircuit.rectangularTraceNorm_sourceOnly_density_sub_le_half_of_count_le}
    \lean{TNLean.PEPS.PairEffect.OriginalCircuit.rectangularTraceNorm_sourceOnly_wholeParty_density_sub_le_half}
    \lean{TNLean.PEPS.PairEffect.sourceGateBudget_eq_of_pos}
    \leanok
    \uses{thm:peps_circuit_effect_error}
    Fix a decomposition of the original output memory into physical and
    discarded factors. For an input $x$ with $\|x\|\le1$, let $\rho$ be
    the original physical density and let $\widetilde\rho$ be the
    replacement density after also discarding its auxiliary registers.
    Then
    \begin{align}
        \|\widetilde\rho-\rho\|_1\le4\delta N.
        \label{eq:peps_source_only_density_error}
    \end{align}
    If $N\le M$ and $\varepsilon>0$, choose
    $\delta=\varepsilon/(8\max\{1,M\})$. The density error is at most
    $\varepsilon/2$. For $M>0$ this is precisely
    $\delta=\varepsilon/(8M)$. The estimate applies in particular to a
    tensor product of normalized vectors, one in each party's entire
    private memory.
\end{theorem}

\begin{proof}\leanok
    \lean{Matrix.appendAuxiliaryVector}
    \lean{Matrix.norm_appendAuxiliaryVector}
    \lean{Matrix.partialTraceRight_appendAuxiliaryVector}
    \lean{Matrix.rectangularTraceNorm_appendAuxiliaryVector_error_le_two}
    \lean{TNLean.PEPS.PairEffect.EffectCircuit.finiteDimensional_auxiliary}
    \lean{TNLean.PEPS.PairEffect.EffectCircuit.originalPhysicalDensity}
    \lean{TNLean.PEPS.PairEffect.EffectCircuit.rectangularTraceNorm_replacement_density_sub_le}
    \lean{TNLean.PEPS.PairEffect.EffectCircuit.replacementPhysicalDensity}
    \lean{TNLean.PEPS.PairEffect.PartyChain.finiteDimensional_sources_prepStack}
    \lean{TNLean.PEPS.PairEffect.auxiliaryCoordinateIso}
    \lean{TNLean.PEPS.PairEffect.auxiliaryCoordinateIso_tmul}
    \lean{TNLean.PEPS.PairEffect.auxiliaryCoordinates}
    \lean{TNLean.PEPS.PairEffect.extendedPhysicalReadout}
    \lean{TNLean.PEPS.PairEffect.extendedPhysicalReadout_append_tmul}
    \lean{TNLean.PEPS.PairEffect.finiteDimensional_mem_append}
    \lean{TNLean.PEPS.PairEffect.finiteDimensional_sources_comp}
    \lean{TNLean.PEPS.PairEffect.finiteDimensional_sources_prepGate}
    \lean{TNLean.PEPS.PairEffect.finiteDimensional_sources_prepWord}
    \lean{TNLean.PEPS.PairEffect.norm_extendedPhysicalReadout_le_one}
    \lean{TNLean.PEPS.PairEffect.rectangularTraceNorm_extendedPhysicalReadout_error_le}
    \lean{TNLean.PEPS.PairEffect.OriginalCircuit.rectangularTraceNorm_sourceOnly_density_sub_le_half}
    \lean{TNLean.PEPS.PairEffect.sourceGateBudget}
    \lean{TNLean.PEPS.PairEffect.sourceGateBudget_pos}
    \lean{TNLean.PEPS.PairEffect.sourceGateBudget_spec}
    \uses{thm:peps_circuit_effect_error}
    Put $u=\widetilde T x$ and $v=\Gamma\otimes Tx$, expressed in the
    common physical and discarded coordinates. Both norms are at most
    one, and (\ref{eq:peps_circuit_effect_error}) gives
    $\|u-v\|\le2\delta N$. For the final partial trace $\operatorname{Tr}_E$,
    \begin{align}
        \|\operatorname{Tr}_E(|u\rangle\langle u|-|v\rangle\langle v|)\|_1
        \le(\|u\|+\|v\|)\|u-v\|\le4\delta N.
        \label{eq:peps_source_only_pure_error}
    \end{align}
    Discarding the normalized factor $\Gamma$ leaves the original
    physical density unchanged. This proves
    (\ref{eq:peps_source_only_density_error}); substituting the chosen
    budget and $N\le M$ gives the final bound. When $M=0$, one has $N=0$
    and the error is zero.
\end{proof}

\begin{definition}[Original local physical dimensions]
    \label{def:peps_family_physical_output}
    \lean{TNLean.PEPS.PairEffect.familyPhysicalLayout}
    \lean{TNLean.PEPS.PairEffect.familyPhysicalOutputLayout}
    \lean{TNLean.PEPS.PairEffect.familyPhysicalListBasis}
    \lean{TNLean.PEPS.PairEffect.familyLabelledPhysicalBasis}
    \leanok
    \uses{def:peps_party_layout}
    Let the original physical space at party $p$ be $H_p=\C^{d_p}$,
    where $d_p$ is a nonnegative integer, and let
    $\mathbf p=(p_1,\ldots,p_m)$ be an ordered list of parties.
    The physical register in position $i$ is $H_{p_i}$, owned by $p_i$.
    Its standard tensor basis has column
    $e_{p_1,z_1}\otimes\cdots\otimes e_{p_m,z_m}\otimes1$
    for position coordinates $z_i\in\{0,\ldots,d_{p_i}-1\}$.
    For an ordering containing each party exactly once, reindex these
    coordinates by assignments $x_p\in\{0,\ldots,d_p-1\}$
    labelled by the original parties, with $z_i=x_{p_i}$.
    Selecting register owners in a region $A$ gives the corresponding
    dependent assignments on $A$.

    Append the actual private registers after the physical registers.
    When their memory is finite-dimensional, with no numerical dimension
    bound, its orthonormal basis and the standard physical tensor basis
    give exact physical-private coordinates. The regional readout retains
    $x|_A$ and discards both $x|_{P\setminus A}$ and these private coordinates.
    All tensor memories retain their terminal scalar factor $\C$.
\end{definition}

\begin{theorem}[Regional dimensions without padding]
    \label{thm:peps_family_regional_dimensions}
    \lean{TNLean.PEPS.PairEffect.OriginalCircuit.rectangularTraceNorm_sourceOnly_family_regional_density_sub_le_half}
    \leanok
    \uses{def:peps_family_physical_output,thm:peps_source_only_density_error}
    The actual physical memory selected by the original party set $A$
    has dimension
    \begin{align}
        \dim\mathcal H_A=\prod_{p\in A}d_p\le D^{|A|}
        \label{eq:peps_family_regional_dimension}
    \end{align}
    whenever $d_p\le D$ on $A$. The regional readout splits the same
    original configuration into its restrictions to $A$ and its
    complement, and discards the actual private memory. With this
    readout, the original whole-party product input and the gate budget
    of Theorem~\ref{thm:peps_source_only_density_error} give replacement
    density error at most $\varepsilon/2$.
\end{theorem}

\begin{proof}\leanok
    \uses{thm:peps_source_only_density_error}
    Selecting owned physical registers filters the original party
    ordering. Its tensor basis is indexed by one choice of $d_p$ values
    at each $p\in A$, so its cardinality is $\prod_{p\in A}d_p$.
    Multiplying the local bounds gives
    (\ref{eq:peps_family_regional_dimension}). Reindexing the full tensor
    basis by the two restrictions of the original configuration gives
    the regional isometry. Apply
    Theorem~\ref{thm:peps_source_only_density_error} to this actual
    partial trace.
\end{proof}

\begin{theorem}[Constant-dimensional regional density estimate]
    \label{thm:peps_constant_family_regional_density}
    \lean{TNLean.PEPS.PairEffect.OriginalCircuit.rectangularTraceNorm_sourceOnly_regional_density_sub_le_half}
    \uses{def:peps_family_physical_output,thm:peps_source_only_density_error}
    Suppose $d_p=d$ at every party and the actual discarded private memory
    is finite-dimensional. For the original whole-party product input,
    contraction circuit and positive error budget of
    Theorem~\ref{thm:peps_source_only_density_error}, the regional readout
    retains the original physical coordinates in $A$ and discards the
    exterior physical and private coordinates. If the original gate count
    is bounded by $M$, the resulting source-only replacement density has
    trace-norm error at most $\varepsilon/2$.
\end{theorem}
\begin{proof}
    \uses{thm:peps_source_only_density_error}
    Specialize the physical dimension family to $d_p=d$, using the same
    complete party ordering and its coverage. Apply the density-error
    estimate to the resulting regional isometry and the original normalized
    whole-party product vector.
\end{proof}

\begin{theorem}[Resources of the constructed source-only circuit]
    \label{thm:peps_source_only_resources}
    \lean{TNLean.PEPS.PairEffect.sum_norm_replacementCoefficients_le}
    \lean{TNLean.PEPS.PairEffect.card_replacementLabels_le}
    \lean{TNLean.PEPS.PairEffect.stackLength_le}
    \lean{TNLean.PEPS.PairEffect.card_replacementLabels_stackLength_le}
    \lean{TNLean.PEPS.PairEffect.OriginalCircuit.isExpansionBounded_replacement}
    \lean{TNLean.PEPS.PairEffect.PreparedSourceGate.length_sources_branchWord_eq_choose}
    \lean{TNLean.PEPS.PairEffect.SourceCircuit.slotCount_eq_choose}
    \lean{TNLean.PEPS.PairEffect.OriginalCircuit.card_sourceLocations_replacement_le}
    \lean{TNLean.PEPS.PairEffect.OriginalCircuit.participationCount_replacement}
    \lean{TNLean.PEPS.PairEffect.OriginalCircuit.participationCount_replacement_le}
    \leanok
    \uses{thm:peps_circuit_effect_error,thm:peps_replacement_participants}
    Suppose each original nonprivate gate has at most $K$ monomials,
    absolute coefficient sum at most $S$, and at most $r$ pair effects
    per monomial. Suppose also that each gate involves at most $b$
    parties and each party participates in at most $b$ nonprivate
    occurrences over the entire circuit.

    With the stack length chosen above, each replacement gate has at
    most $Km^r$ monomials and absolute coefficient sum at most $S$.
    Moreover,
    \begin{align}
        m    & \le(rS/\delta)^2+2,                 \\
        Km^r & \le K\bigl((rS/\delta)^2+2\bigr)^r.
        \label{eq:peps_replacement_resource_bound}
    \end{align}
    After combining sources on equal party pairs, a gate with $q$
    participating parties has exactly $\binom q2$ common source slots,
    including one-dimensional padding. Thus the entire replacement
    has at most $N\binom b2$ source positions. Each party participates
    in exactly as many expanded gate occurrences as originally, so
    its bound $b$ over the entire circuit is preserved. At each actual
    chronological occurrence, the branch cardinality and coefficient sum
    satisfy the same bounds. The participation count is the number of
    incident nonprivate occurrences, with repeated uses counted separately.
\end{theorem}

\begin{proof}\leanok
    \lean{TNLean.PEPS.PairEffect.exists_prepared_effectReplacement_uniform_resources}
    \lean{TNLean.PEPS.PairEffect.length_wordList_le}
    \lean{TNLean.PEPS.PairEffect.sum_norm_coeff_wordList}
    \lean{TNLean.PEPS.PairEffect.sum_norm_replacementCoefficients}
    \lean{TNLean.PEPS.PairEffect.OriginalCircuit.IsMonomialBounded}
    \lean{TNLean.PEPS.PairEffect.SourceCircuit.IsExpansionBounded}
    \lean{TNLean.PEPS.PairEffect.SourceCircuit.isExpansionBounded_castLayouts}
    \lean{TNLean.PEPS.PairEffect.SourceCircuit.isExpansionBounded_exchangeBlocks}
    \lean{TNLean.PEPS.PairEffect.SourceCircuit.isExpansionBounded_frameList}
    \lean{TNLean.PEPS.PairEffect.SourceCircuit.isExpansionBounded_ofLocalWord}
    \lean{TNLean.PEPS.PairEffect.OriginalCircuit.card_nonprivateLocations}
    \lean{TNLean.PEPS.PairEffect.OriginalCircuit.natCard_nonprivateLocations}
    \lean{TNLean.PEPS.PairEffect.OriginalCircuit.nonprivateLocationsFintype}
    \lean{TNLean.PEPS.PairEffect.OriginalCircuit.participationCount}
    \lean{TNLean.PEPS.PairEffect.OriginalCircuit.participationLocationsEquiv}
    \lean{TNLean.PEPS.PairEffect.SourceCircuit.participationCount}
    \lean{TNLean.PEPS.PairEffect.PreparedSourceGate.length_slots}
    \lean{TNLean.PEPS.PairEffect.PreparedSourceGate.length_sources_branchWord}
    \lean{TNLean.PEPS.PairEffect.PreparedSourceGate.length_sources_branchWord_le_choose}
    \lean{TNLean.PEPS.PairEffect.SourceCircuit.card_sourceLocations_eq_sum_choose}
    \lean{TNLean.PEPS.PairEffect.SourceCircuit.card_sourceLocations_le}
    \lean{TNLean.PEPS.PairEffect.SourceCircuit.IsExpansionBounded.at_gate}
    \lean{TNLean.PEPS.PairEffect.SourceCircuit.participationCount_eq_card_incidentGates}
    \uses{cor:peps_pair_effects_rescaling,thm:peps_replacement_participants}
    Expanding the permutation averages preserves the absolute sum of
    the original coefficients. The common rescaling has modulus
    $(1+\delta)^{-1}\le1$, so it cannot increase that sum. The monomial
    count is at most $Km^r$ by the gate expansion. The ceiling estimate
    gives (\ref{eq:peps_replacement_resource_bound}). Inducting on the
    original chronology carries both bounds through composition and
    retained spectators.

    Each prepared gate has one common slot for every unordered pair of
    distinct participating parties. Summing $\binom q2\le\binom b2$
    over its $N$ occurrences bounds all source positions. The occurrence
    correspondence of Theorem~\ref{thm:peps_replacement_participants}
    restricts to a bijection on the occurrences involving any fixed
    party, proving exact preservation of its count.
\end{proof}

% These statements establish the effect-elimination part of the proof of
% Theorem 5.2. They do not establish the subsequent sampling construction,
% uniform polynomial bounds, or the full distributed density-compression theorem.
